\documentclass[12pt]{article}
\usepackage{amsmath, amssymb}
\usepackage[margin=1in]{geometry}
\setlength{\parskip}{6pt}
\title{QUBO Formulation: Network Community Detection}
\date{}
\begin{document}
\maketitle

\subsection*{Network Community Detection}

\paragraph{Problem Statement.}
Given a graph and a set of $K$ candidate community partitions, each associated with a precomputed modularity-like quality score $S_k$, the task is to select exactly one partition that maximizes the overall community quality. The selection must satisfy the requirement that precisely one candidate is chosen, and the objective is to return the optimal partition along with its maximum score.

\paragraph{Decision Variables.}
For each candidate partition $k \in \{0, 1, \dots, K-1\}$, introduce a binary decision variable $y_k \in \{0,1\}$. The variable $y_k$ equals $1$ if candidate partition $k$ is selected, and $0$ otherwise. The vector of decision variables is denoted $\mathbf{y} = (y_0, y_1, \dots, y_{K-1})^\top$.

\paragraph{QUBO Derivation.}
\textbf{Step 1 -- Objective Function.} The original goal is to maximize the total modularity-like score:
\begin{align}
    \max_{\mathbf{y} \in \{0,1\}^K} \sum_{k=0}^{K-1} S_k y_k.
\end{align}
Since QUBO solvers minimize energy, we negate the objective:
\begin{align}
    \min_{\mathbf{y} \in \{0,1\}^K} -\sum_{k=0}^{K-1} S_k y_k.
\end{align}

\textbf{Step 2 -- Constraints.} The requirement that exactly one partition is selected is expressed as:
\begin{align}
    \sum_{k=0}^{K-1} y_k = 1.
\end{align}

\textbf{Step 3 -- Penalty Term Expansion.} We enforce the constraint using a quadratic penalty with weight $\lambda > 0$:
\begin{align}
    P(\mathbf{y}) = \lambda \left( \sum_{k=0}^{K-1} y_k - 1 \right)^2.
\end{align}
Expanding the square yields:
\begin{align}
    P(\mathbf{y}) = \lambda \left( \sum_{k=0}^{K-1} y_k^2 + \sum_{k \neq j} y_k y_j + 1 - 2\sum_{k=0}^{K-1} y_k \right).
\end{align}
Using the idempotent property of binary variables, $y_k^2 = y_k$, the expression simplifies to:
\begin{align}
    P(\mathbf{y}) = \lambda \left( 1 - \sum_{k=0}^{K-1} y_k + \sum_{k \neq j} y_k y_j \right).
\end{align}
Noting that $\sum_{k \neq j} y_k y_j = 2 \sum_{0 \le k < j \le K-1} y_k y_j$, the penalty contribution becomes:
\begin{align}
    P(\mathbf{y}) = \lambda - \lambda \sum_{k=0}^{K-1} y_k + 2\lambda \sum_{0 \le k < j \le K-1} y_k y_j.
\end{align}

\textbf{Step 4 -- Combined Hamiltonian.} Adding the negated objective and the penalty term gives the QUBO Hamiltonian $H(\mathbf{y})$:
\begin{align}
    H(\mathbf{y}) = -\sum_{k=0}^{K-1} S_k y_k - \lambda \sum_{k=0}^{K-1} y_k + 2\lambda \sum_{0 \le k < j \le K-1} y_k y_j + \lambda.
\end{align}

\textbf{Step 5 -- Q Matrix and Offset.} Reading off the coefficients from the general form $H(\mathbf{y}) = \mathbf{y}^\top Q \mathbf{y} + c$, the upper-triangular Q matrix entries and constant offset are:
\begin{align}
    Q_{k,k} &= -S_k - \lambda, & \forall k \in \{0, \dots, K-1\}, \\
    Q_{k,j} &= 2\lambda, & \forall 0 \le k < j \le K-1, \\
    c &= \lambda.
\end{align}

\paragraph{Penalty Weight Justification.}
The penalty weight must be chosen as $\lambda > \max_{k} |S_k|$. This condition ensures that the energy penalty incurred by any infeasible assignment strictly exceeds the maximum possible gain achievable by the objective function. Specifically, if zero variables are selected, the penalty contributes $+\lambda$ to the energy. If two or more variables are selected, the quadratic term $2\lambda \sum_{k<j} y_k y_j$ contributes at least $2\lambda$, which dominates any linear objective term bounded by $\max_k |S_k|$. Thus, $\lambda > \max_k |S_k|$ guarantees that feasible solutions always have lower energy than infeasible ones.

\paragraph{Correctness Argument.}
Minimizing $H(\mathbf{y})$ over $\{0,1\}^K$ recovers the optimal partition because the penalty term is zero if and only if $\sum_k y_k = 1$, meaning all feasible solutions have unpenalized energy equal to $-\sum_k S_k y_k$. Infeasible solutions incur a penalty of at least $\lambda$, which by construction exceeds the maximum possible reduction in the objective term. Consequently, the global minimum of $H(\mathbf{y})$ must occur at a feasible point, and among feasible points, the minimum corresponds to the partition with the highest modularity-like score.

\paragraph{Illustrative Example.}
Consider an instance with $K=3$ candidate partitions and scores $S_0 = 0.4$, $S_1 = 0.7$, and $S_2 = 0.5$. We choose $\lambda = 1.0$, which satisfies $\lambda > \max_k |S_k|$. Substituting these values into the general formulas yields the concrete Hamiltonian:
\begin{align}
    H(\mathbf{y}) &= -0.4 y_0 - 0.7 y_1 - 0.5 y_2 - 1.0 y_0 - 1.0 y_1 - 1.0 y_2 \nonumber \\
    &\quad + 2.0 y_0 y_1 + 2.0 y_0 y_2 + 2.0 y_1 y_2 + 1.0.
\end{align}
The resulting Q matrix entries and offset are:
\begin{align}
    Q_{0,0} &= -1.4, \quad Q_{1,1} = -1.7, \quad Q_{2,2} = -1.5, \\
    Q_{0,1} &= 2.0, \quad Q_{0,2} = 2.0, \quad Q_{1,2} = 2.0, \\
    c &= 1.0.
\end{align}
Evaluating $H(\mathbf{y})$ over the feasible subspace $\{ \mathbf{y} \mid \sum_k y_k = 1 \}$ gives energies $H(1,0,0) = -0.4$, $H(0,1,0) = -0.7$, and $H(0,0,1) = -0.5$. The global minimum is $-0.7$ at $\mathbf{y}^* = (0,1,0)^\top$, correctly selecting candidate partition $1$ as the optimal community assignment.
\end{document}